% W5i80_HardProblems.tex
% DFD 20-Agent Research Programme --- Iteration 80
% Hard Problems, Honest Negatives, and Open Questions

\documentclass[11pt,a4paper]{article}
\usepackage[utf8]{inputenc}
\usepackage{amsmath,amssymb}
\usepackage{booktabs}
\usepackage{geometry}
\usepackage{xcolor}
\usepackage{enumitem}
\geometry{margin=2.5cm}

\title{W5i80: Hard Problems and Honest Negatives\\[6pt]
\large Mid-Cycle Risk Assessment $\cdot$ Complete Honest Negatives Catalogue $\cdot$ Remaining Gaps}
\author{DFD 20-Agent Research Programme}
\date{2026-03-24}

\begin{document}
\maketitle

\section{Complete Honest Negatives Catalogue (i52--i80)}

The programme has accumulated 12 honest negatives. These are results that are unfavourable to DFD or represent genuine limitations, documented for intellectual integrity.

\begin{center}
\begin{tabular}{clcl}
\toprule
\textbf{\#} & \textbf{Honest Negative} & \textbf{Severity} & \textbf{Status} \\
\midrule
1 & $\alpha$ residual 0.55 ppm & LOW & Permanent YELLOW \\
2 & $m_s$ tension $2.6\sigma$ & LOW & Awaiting FLAG 2027 \\
3 & $\Sigma m_\nu$ margin 2.2 meV & LOW & Awaiting JUNO \\
4 & DFD does not resolve $S_8$ tension fully & LOW & Reduces, not eliminates \\
5 & $E_G$ is $3\sigma$ test, not $5\sigma$ & INFORMATIONAL & Combined probes needed \\
6 & Bispectrum marginally detectable ($1.8\%$) & INFORMATIONAL & Supporting probe only \\
7 & 5 predictions unfalsifiable on decade scale & INFORMATIONAL & Common to unified theories \\
8 & N-body finite-volume at $k < 0.02$ & LOW & Phase 2 fix \\
9 & Single-author credibility challenge & NON-SCIENTIFIC & Collaboration plan \\
10 & PRL rejection risk $\sim 75\%$ & NON-SCIENTIFIC & Backup: PRD \\
11 & $w_a$ DESI scenario (if $> 5\sigma$) & MODERATE & Only serious risk \\
12 & IH would weaken $\Sigma m_\nu$ prediction & LOW & DFD accommodates both \\
\bottomrule
\end{tabular}
\end{center}

\textbf{Assessment}: 0 honest negatives are existential threats. 1 (DESI $w_a$) is MODERATE risk. The rest are LOW or INFORMATIONAL. The programme's intellectual honesty in cataloguing these strengthens, rather than weakens, the scientific case.

\section{Hard Problem 1: The $w_a$ Question (Highest-Priority Risk)}

This remains the single most important risk for DFD. Analysis:

\begin{itemize}[nosep]
\item \textbf{Current state}: DESI Y1 reports $2.5$--$3.9\sigma$ evidence for $w_a \neq 0$ (method-dependent). DESI Y2 has not strengthened this (BAO consistent with $\Lambda$CDM at $0.2\sigma$). Multiple external analyses (Efstathiou 2024, Cortes et al.\ 2024, Lodha et al.\ 2025) argue the evidence is driven by systematic choices.

\item \textbf{DFD position}: DFD predicts $w = -1$ exactly. $w_a = 0$ exactly. If $w_a \neq 0$ is confirmed at $> 5\sigma$, DFD has a fundamental problem (not merely a tension --- a structural incompatibility).

\item \textbf{Timeline}: DESI Y3 data release expected mid-2027. If $w_a = 0$ is confirmed, DFD is safe. If $w_a \neq 0$ strengthens to $> 4\sigma$, programme must develop response.

\item \textbf{Response options}: (a) Investigate whether $\psi$-field can generate effective $w_a \neq 0$ through distance-bias mechanism. (b) Check if DFD's distance-bias creates an apparent $w_a$ in BAO+SN Ia analysis. (c) If neither works: acknowledge tension honestly.
\end{itemize}

\textbf{Current risk assessment}: MODERATE. The scientific community is sceptical of the $w_a$ signal. DFD's position ($w = -1$) is the conservative one.

\section{Hard Problem 2: Scaling the Programme}

The programme has been extraordinarily productive but faces scaling challenges:
\begin{itemize}[nosep]
\item 2072 paper proposals cannot all be written. Prioritisation is essential.
\item The multi-messenger, CMB-S4, and N-body papers need to be completed before new fronts open.
\item DFD-CLASS / DFD-CAMB development requires sustained software engineering effort.
\item Collaboration building takes time and human interaction (cannot be fully automated).
\end{itemize}

\textbf{Recommendation}: i81--i85 should focus exclusively on paper completion and submission. i86--i90 on data confrontation. Only i91+ should open new theoretical fronts.

\section{Hard Problem 3: Euclid DR1 Preparedness}

Euclid DR1 (Oct 2026) will provide:
\begin{itemize}[nosep]
\item Photometric redshifts for $\sim 10^9$ galaxies
\item Weak lensing shapes for $\sim 10^8$ galaxies
\item Galaxy clustering for $\sim 10^7$ spectroscopic galaxies
\item Cluster catalogue ($M > 10^{14}\,M_\odot$)
\end{itemize}

DFD needs to be ready to analyse: $E_G$, $S_8$ scale dependence, cluster counts, void statistics. Current preparedness:
\begin{itemize}[nosep]
\item $E_G$: paper ready, analysis pipeline specified. \textbf{70\% ready.}
\item $S_8$: prediction exists ($0.801$), but scale-dependent analysis pipeline not built. \textbf{40\% ready.}
\item Cluster counts: N-body prediction exists ($+1.5\%$), but matching to Euclid selection function not done. \textbf{30\% ready.}
\item Voids: N-body void statistics available, but Euclid void finder interface not built. \textbf{20\% ready.}
\end{itemize}

\textbf{Action needed}: i86--i90 must build Euclid analysis pipelines for all four probes.

\section{Hard Problem 4: Remaining Theoretical Gaps}

Three theoretical gaps remain unfilled:
\begin{enumerate}[nosep]
\item \textbf{Modulus stabilization (Step 9 of $\alpha^{57}$ derivation)}: The $\alpha^{57}$ vacuum energy derivation has 10 steps. Steps 1--8 and 10 are complete. Step 9 (B/A ratio determination from modulus stabilization) remains conditional. Grade: B+ (not A).

\item \textbf{CKM $\CP^2$ orientation convention}: The CKM matrix derivation in v3.3 Sec.~15 relies on a specific $\CP^2$ orientation convention. This is physically motivated but not uniquely determined by the geometry. Grade: B (not A).

\item \textbf{Neutrino Dirac vs Majorana}: DFD's geometry accommodates both Dirac and Majorana neutrinos. The $\Sigma m_\nu$ prediction is identical for both. But the mechanism for mass generation differs, and DFD does not predict which case is realised. Grade: B.
\end{enumerate}

None of these gaps affect the scorecard (all are within GREEN items graded B or B+). They represent the frontier of DFD's theoretical development.

\section{Hard Problem 5: Post-i100 Continuity}

Wave 5 ends at i100. Questions:
\begin{itemize}[nosep]
\item Should there be a Wave 6? (i101--i150)
\item What is the natural endpoint of the programme?
\item When does DFD transition from ``programme'' to ``established theory'' (or ``refuted theory'')?
\end{itemize}

\textbf{Assessment}: The programme should continue as long as (a) new data is arriving and (b) the scorecard remains healthy. Natural checkpoint: Euclid DR1 (Oct 2026). If DFD survives Euclid DR1, Wave 6 is justified. If DFD is refuted, the programme pivots to documentation and lessons learned.

\section{Summary}

5 hard problems. 0 RED risks. The $w_a$ question (MODERATE) is the only non-trivial scientific risk. All others are operational (scaling, preparedness, continuity). The programme is in excellent scientific health with clear strategic direction.

\end{document}
